\documentclass[11pt,a4paper]{article}

% ----------------------------------------------------------------------------
%  Research Story: Evolution of a Crypto Trading Strategy
%  SG vs HK Hackathon 2026 -- Strategy 1
%
%  Compile with:  pdflatex research_story.tex   (run twice for the ToC/refs)
% ----------------------------------------------------------------------------

\usepackage[margin=1in]{geometry}
\usepackage{amsmath,amssymb}
\usepackage{booktabs}
\usepackage{xcolor}
\usepackage{graphicx}
\usepackage{enumitem}
\usepackage{listings}
\usepackage[hidelinks]{hyperref}

\definecolor{codebg}{rgb}{0.96,0.96,0.96}
\definecolor{accent}{rgb}{0.10,0.30,0.55}

\lstset{
  basicstyle=\ttfamily\footnotesize,
  backgroundcolor=\color{codebg},
  breaklines=true,
  frame=single,
  framesep=4pt,
  columns=fullflexible,
  showstringspaces=false,
}

\newcommand{\param}[1]{\texttt{#1}}
\newcommand{\code}[1]{\texttt{#1}}

\title{\textbf{Don't Predict, React:}\\[2pt]
       \large The Research Evolution of a Regime-Aware\\ Crypto Momentum Strategy}
\author{Strategy 1 --- SG vs HK Hackathon 2026\\[2pt]
        \small \texttt{<Your Name / Team>}}
\date{March 2026}

\begin{document}
\maketitle

\begin{abstract}
\noindent
This report documents the full research journey behind ``Strategy 1,'' an
algorithmic trading bot built for the Roostoo mock-exchange during the SG~vs~HK
2026 hackathon. The strategy was not designed in one shot; it emerged through
five distinct phases, each one forced by a concrete, observed failure of its
predecessor. We began with the most intuitive prior --- mean reversion --- and
progressively discovered that the edge in crypto markets is fundamentally
\emph{regime-dependent}. The path runs: (1)~mean reversion, (2)~mean reversion
with rolling linear-regression detrending, (3)~an inversion to momentum via
volatility breakout, (4)~a regime-switching ensemble that unifies both views,
and finally (5)~a deliberate simplification to gated breakout-momentum hardened
with institutional-grade risk controls. We present the hypothesis, mechanics,
mathematics, and failure mode of each phase, and conclude with the final
production specification.
\end{abstract}

\tableofcontents
\bigskip

% ============================================================================
\section{Introduction}
% ============================================================================

\subsection{The problem}
The objective was to build a fully automated trading bot that grows a
\$50{,}000 paper portfolio on the Roostoo mock exchange, trading a small
universe of liquid spot pairs (\code{BTC}, \code{ETH}, \code{SOL}, \code{BNB})
against \code{USD}. The bot polls live prices on a $\sim$30-second cycle, places
\code{LIMIT} orders through a signed REST API, and must survive unattended for
extended periods without human intervention.

This imposes three constraints that shaped every design decision:
\begin{enumerate}[noitemsep]
  \item \textbf{No look-ahead.} Decisions at tick $t$ may only use information
        available up to and including $t$.
  \item \textbf{Capital preservation first.} A blown-up bot scores zero. Drawdown
        control is not optional polish --- it is a primary objective.
  \item \textbf{Robustness over cleverness.} The bot trades unsupervised, so every
        edge case (stale prices, partial fills, API errors) must be handled
        mechanically.
\end{enumerate}

\subsection{Research philosophy}
Rather than search for a single ``magic'' indicator, we treated the project as a
sequence of falsifiable hypotheses. Each strategy was a bet about \emph{how crypto
prices behave}; we deployed it, observed where it broke, and let that failure
dictate the next hypothesis. The result is a narrative of
\emph{thesis $\rightarrow$ antithesis $\rightarrow$ synthesis $\rightarrow$
refinement}, summarised in Table~\ref{tab:journey}.

\begin{table}[h]
\centering
\small
\begin{tabular}{@{}llll@{}}
\toprule
\textbf{Phase} & \textbf{Hypothesis} & \textbf{Mechanism} & \textbf{Why we moved on} \\
\midrule
1. Mean reversion       & Prices snap back to a mean      & $Z$-score on rolling mean        & Trends run it over \\
2. MR + linear reg.     & The mean \emph{drifts}          & $Z$-score on regression residual & Still structurally short trend \\
3. Momentum (breakout)  & Ride proven moves               & Donchian + trailing ATR stop     & Wins in trends, idle in chop \\
4. Regime ensemble      & Pick the right tool per regime  & BB-width switch: MR $\leftrightarrow$ momentum & Complex; MR leg was weak \\
5. Gated momentum       & Momentum, but only when it pays & Breakout + regime on/off gate    & \textbf{Final} \\
\bottomrule
\end{tabular}
\caption{The five-phase research journey for Strategy 1.}
\label{tab:journey}
\end{table}

\noindent
\textit{Scope note.} This document covers \textbf{Strategy 1} only --- the
breakout/momentum lineage authored for this project. A separate
\code{strategy 2/} module (an EWMA high-frequency-momentum signal) was contributed
independently by a teammate and is outside the scope of this narrative; it is
discussed briefly in Appendix~\ref{app:strat2} for completeness.

% ============================================================================
\section{Phase 1 --- Mean Reversion: The Naive Prior}
% ============================================================================

\subsection{Hypothesis}
The most intuitive starting point: over short horizons, crypto prices
\emph{overreact} and then revert toward a local fair value. If so, we should buy
when price is unusually low relative to its recent average and sell when it is
unusually high.

\subsection{Mechanics}
Let $P_t$ be the price and $\mu_t,\sigma_t$ the rolling mean and standard
deviation over a window $W$:
\begin{equation}
  \mu_t = \frac{1}{W}\sum_{i=0}^{W-1} P_{t-i}, \qquad
  \sigma_t = \sqrt{\frac{1}{W}\sum_{i=0}^{W-1}\bigl(P_{t-i}-\mu_t\bigr)^2}.
\end{equation}
We then form the standardized deviation (a $Z$-score) and trade its extremes:
\begin{equation}
  Z_t = \frac{P_t - \mu_t}{\sigma_t},\qquad
  \text{signal}_t =
  \begin{cases}
    \text{BUY}  & Z_t < -z^{*}\\
    \text{SELL} & Z_t > +z^{*}\\
    \text{HOLD} & \text{otherwise.}
  \end{cases}
\end{equation}

\subsection{What we learned}
Crypto \textbf{trends hard}. A flat rolling mean lags a directional move, so the
$Z$-score stays negative all the way down a sell-off --- the strategy keeps issuing
BUY signals into a falling market, ``catching a falling knife'' and averaging into
losers. Mean reversion implicitly assumes a stationary mean; crypto violates that
assumption precisely when it matters most. This failure motivated Phase~2: if the
mean is the problem, model the mean better.

% ============================================================================
\section{Phase 2 --- Mean Reversion with Rolling Linear Regression}
% ============================================================================

\subsection{Hypothesis}
The Phase~1 failure had a clear diagnosis: the ``mean'' is not flat, it
\emph{drifts}. By fitting a local trend line and measuring deviations from
\emph{that} line (rather than from a flat average), we should be able to keep the
mean-reversion logic while accounting for short-term drift.

\subsection{Mechanics}
Over a lookback of $L=60$ bars we fit an ordinary least-squares trend line,
regressing price on time $\tau$:
\begin{equation}
  \hat{P}_\tau = \hat\alpha + \hat\beta\,\tau,
  \qquad
  (\hat\alpha,\hat\beta) = \arg\min_{\alpha,\beta}\sum_{\tau=t-L+1}^{t}\bigl(P_\tau-\alpha-\beta\tau\bigr)^2 .
\end{equation}
The \emph{residual} is the part of price the local trend does \emph{not} explain:
\begin{equation}
  r_t = P_t - \hat{P}_t,\qquad
  Z_t = \frac{r_t}{\sigma_r},\quad
  \sigma_r = \operatorname{std}\bigl(r_{t-L+1:t}\bigr).
\end{equation}
We trade the residual $Z$-score and, crucially, add \textbf{position management}:
the size scales with the magnitude of the dislocation, so we add more as the
price strays further from the regression line.

\subsection{What we learned}
Detrending helped at the margin but did not fix the core defect. Mean reversion is
\emph{structurally short trend}: it profits only if the dislocation closes. In a
genuine directional regime the residual keeps widening in the same direction, and
the ``add more as it strays'' rule amplifies the loss. The decisive insight emerged
here: \textbf{the edge is regime-dependent.} Fading moves works in ranging markets
and is actively dangerous in trending ones. If fighting trends loses money, perhaps
we should do the opposite.

% ============================================================================
\section{Phase 3 --- Momentum via Volatility Breakout: The Inversion}
% ============================================================================

\subsection{Hypothesis}
We flipped the thesis. Instead of betting that moves reverse, bet that
\emph{proven} moves \emph{continue}. Do nothing until the market commits to a
direction; then ride it with a stop that lets winners run and cuts losers
mechanically.

\subsection{Mechanics}
\paragraph{Donchian channel.} Define the entry and exit channels from the
trailing highs/lows (excluding the current bar) over lookbacks
$\param{ENTRY\_LB}=120$ and $\param{EXIT\_LB}=48$:
\begin{equation}
  H_t = \max\bigl(P_{t-\param{ENTRY\_LB}}, \dots, P_{t-1}\bigr),\qquad
  L_t = \min\bigl(P_{t-\param{EXIT\_LB}}, \dots, P_{t-1}\bigr).
\end{equation}
\paragraph{Volatility (ATR proxy).} A mean-absolute-change estimate over
$\param{ATR\_WIN}=20$ bars:
\begin{equation}
  \mathrm{ATR}_t = \frac{1}{\param{ATR\_WIN}}\sum_{i=1}^{\param{ATR\_WIN}}\bigl|P_{t-i+1}-P_{t-i}\bigr|.
\end{equation}
\paragraph{Entry.} Go long when price breaks the upper channel, but only after
three confirmations stack up:
\begin{align}
  &\text{(i) breakout:}    && P_t > H_t \text{ for } \param{CONFIRM\_TICKS}=2 \text{ consecutive ticks};\\
  &\text{(ii) strength:}   && \frac{P_t - H_t}{\mathrm{ATR}_t} \ge \param{MIN\_BREAKOUT\_STR}=0.5;\\
  &\text{(iii) vol gate:}  && \frac{\mathrm{ATR}_t}{P_t} \ge \param{MIN\_VOL}.
\end{align}
\paragraph{Exit.} A position is closed at the highest of a trailing stop or the
exit channel, or at a profit target:
\begin{equation}
  \text{stop}_t = \text{peak} - \param{STOP\_MULT}\cdot\mathrm{ATR}_t,\quad
  \text{target}_t = P_{\text{entry}} + \param{PROFIT\_TARGET}\cdot\mathrm{ATR}_t,
\end{equation}
with $\param{STOP\_MULT}=3.0$ and $\param{PROFIT\_TARGET}=2.0$.

\subsection{Why this is elegant}
In flat or choppy markets, price never clears the channel: \textbf{zero trades,
zero loss}. The breakout strategy is therefore \emph{naturally defensive in
exactly the chop regime where mean reversion also struggled} --- but for the
opposite reason. We had now covered the two opposing regimes with two opposing
tools.

% ============================================================================
\section{Phase 4 --- Regime Conditioning: The Synthesis}
% ============================================================================

\subsection{Hypothesis}
We now possessed two strategies, each winning in one regime and losing in the
other: mean reversion in ranges, momentum in trends. The natural synthesis is a
system that \emph{detects} the prevailing regime and deploys the appropriate tool.

\subsection{Mechanics}
\paragraph{Regime detector.} We use Bollinger-band width --- the spread of the
bands relative to their mean --- as a trend/range classifier over
$\param{REGIME\_WIN}=20$ bars:
\begin{equation}
  \mathrm{BBwidth}_t = \frac{4\,\sigma_t}{\mu_t},\qquad
  \text{regime} =
  \begin{cases}
    \text{TREND} & \mathrm{BBwidth}_t > \param{BREAKOUT\_THRESH}=0.025\\
    \text{RANGE} & \text{otherwise.}
  \end{cases}
\end{equation}
The raw width is smoothed by an EMA to prevent the bot from thrashing between
regimes on noise.
\paragraph{Trending branch (momentum).} EMA crossover $\mathrm{EMA}_5 >
\mathrm{EMA}_{20}$ sets direction, gated by an RSI overbought filter and a
``BTC must also be trending up'' market filter.
\paragraph{Ranging branch (mean reversion).} Bollinger-band position sizing,
$w = (\text{upper}-P)/(\text{upper}-\text{lower})$, with an RSI confluence
requirement near the lower band and a minimum-bandwidth gate so we only fade when
the band is wide enough to cover costs.
\paragraph{Risk overlay.} Volatility targeting, drawdown-based de-sizing, per-coin
exposure caps, and trade cooldowns sit on top of both branches.

\subsection{What we learned}
The ensemble worked, but it carried two costs: \emph{complexity} (two full
strategies, a switch, and many parameters --- a real overfitting risk on limited
data) and the empirical observation that the \emph{mean-reversion leg was the
weaker half}. The momentum branch and the regime gate were doing most of the work.

% ============================================================================
\section{Phase 5 --- Gated Momentum: Refinement and Hardening}
% ============================================================================

\subsection{Hypothesis}
Having \emph{earned} the complexity of the full ensemble, we earned the right to
cut it back. Keep the part that works --- breakout momentum --- and demote the
regime detector from a \emph{switch} to an \emph{on/off gate}: trade momentum only
when the volatility regime confirms a real trend; otherwise sit in cash.

\subsection{Mechanics}
\paragraph{Regime gate.} Average the per-coin volatility across the universe and
smooth it over $\param{REGIME\_WINDOW}=20$ ticks:
\begin{equation}
  \overline{v}_t = \frac{1}{N}\sum_{c=1}^{N}\frac{\mathrm{ATR}^{(c)}_t}{P^{(c)}_t},\qquad
  \text{regime} =
  \begin{cases}
    \text{ACTIVE} & \widetilde{v}_t \ge \param{REGIME\_VOL\_THRESH}=0.001\\
    \text{FLAT}   & \text{otherwise (no new entries).}
  \end{cases}
\end{equation}
\paragraph{Risk-parity sizing.} Every trade risks a fixed fraction of equity. With
$\param{RISK\_PER\_TRADE}=1\%$ and the stop distance $\param{STOP\_MULT}\cdot
\mathrm{ATR}$, the quantity is
\begin{equation}
  q = \min\!\left(
        \frac{E \cdot \param{RISK\_PER\_TRADE}}{\param{STOP\_MULT}\cdot \mathrm{ATR}_t},\;
        \frac{E \cdot \param{MAX\_ALLOC}}{P_t}
      \right),
\end{equation}
where $E$ is current equity and $\param{MAX\_ALLOC}=24\%$ caps per-coin exposure.
\paragraph{Cross-sectional selection.} Among confirmed breakouts, rank by breakout
strength and hold up to $\param{TOP\_K}=4$ positions, preferring the strongest.

\subsection{The final architecture}
This is the production bot (\code{main.py}). It is, in one sentence:
\emph{a cross-sectional Donchian breakout system that only trades in confirmed
high-volatility regimes, sizes by risk parity, and exits via trailing ATR stops,
channel breaks, or profit targets --- wrapped in circuit breakers.}

% ============================================================================
\section{Final Strategy Specification}
% ============================================================================

\subsection{Signal and execution pipeline}
Each $\sim$30s tick the bot:
\begin{enumerate}[noitemsep]
  \item cancels stale orders and fetches fresh prices (with retry/backoff);
  \item runs stale-data detection (skips trading if prices are frozen);
  \item updates price history, ATR, Donchian channels, and trailing peaks;
  \item \textbf{checks exits first} (stops/channel/target) --- always, even when entries are halted;
  \item evaluates the regime gate and circuit breakers;
  \item checks confirmed-breakout entries and sizes them by risk parity;
  \item reconciles pending fills against the exchange balance and reports equity/P\&L.
\end{enumerate}

\subsection{Parameters}
\begin{table}[h]
\centering
\small
\begin{tabular}{@{}llp{7.2cm}@{}}
\toprule
\textbf{Parameter} & \textbf{Value} & \textbf{Role} \\
\midrule
\param{ENTRY\_LB}           & 120     & Donchian entry-channel lookback \\
\param{EXIT\_LB}            & 48      & Donchian exit-channel lookback (tighter) \\
\param{ATR\_WIN}            & 20      & Volatility (ATR) window \\
\param{TOP\_K}              & 4       & Max simultaneous positions \\
\param{RISK\_PER\_TRADE}    & 1\%     & Equity risked per trade \\
\param{MAX\_ALLOC}          & 24\%    & Per-coin exposure cap \\
\param{STOP\_MULT}          & 3.0     & Trailing stop $=$ peak $-$ STOP\_MULT$\times$ATR \\
\param{PROFIT\_TARGET}      & 2.0     & Take profit at entry $+$ 2$\times$ATR \\
\param{MIN\_BREAKOUT\_STR}  & 0.5     & Min breakout strength (ATR units) \\
\param{CONFIRM\_TICKS}      & 2       & Ticks above channel before entry \\
\param{COOLDOWN}            & 60      & Ticks after stop-out before re-entry \\
\param{REGIME\_VOL\_THRESH} & 0.001   & Avg vol needed for ACTIVE regime \\
\param{REGIME\_WINDOW}      & 20      & Smoothing window for regime vol \\
\bottomrule
\end{tabular}
\caption{Final production parameters (\code{main.py}).}
\label{tab:params}
\end{table}

\subsection{Risk management}
Capital preservation is enforced by a layered set of mechanical controls:
\begin{itemize}[noitemsep]
  \item \textbf{Circuit breakers:} halt new entries on $>5\%$ daily loss,
        $>10\%$ drawdown from peak, or $5$ consecutive losing trades.
  \item \textbf{Trailing ATR stop \& profit target:} bound each trade's downside
        and lock in gains.
  \item \textbf{Risk-parity sizing:} every trade risks the same $1\%$ of equity.
  \item \textbf{Re-entry cooldown:} prevents whipsaw re-entry after a stop-out.
  \item \textbf{Operational safety:} pending-fill reconciliation, stale-price
        detection, dust handling, API retry with exponential backoff, and a
        graceful-shutdown handler with final state logging.
\end{itemize}

% ============================================================================
\section{Results}
% ============================================================================
\noindent
\textit{Populate this section with your own backtest output. The narrative above
is the skeleton; these numbers are the proof. Suggested table:}

\begin{table}[h]
\centering
\small
\begin{tabular}{@{}lccc@{}}
\toprule
\textbf{Phase / Strategy} & \textbf{Total Return} & \textbf{Sharpe} & \textbf{Max Drawdown} \\
\midrule
1. Mean reversion              & \texttt{<--->} & \texttt{<--->} & \texttt{<--->} \\
2. MR + linear regression      & \texttt{<--->} & \texttt{<--->} & \texttt{<--->} \\
3. Volatility breakout         & \texttt{<--->} & \texttt{<--->} & \texttt{<--->} \\
4. Regime ensemble             & \texttt{<--->} & \texttt{<--->} & \texttt{<--->} \\
5. Gated momentum (final)      & \texttt{<--->} & \texttt{<--->} & \texttt{<--->} \\
\bottomrule
\end{tabular}
\caption{Backtest comparison across phases (to be filled in).}
\end{table}

\noindent
For the breakout family, \code{Backtesting.py} provides a walk-forward harness on
historical 1-minute data; reported figures in development included Sharpe ratios of
$+18.76$ (7d), $+5.80$ (1mo), and $+1.75$ (60d) windows. \emph{Verify and update
these against your final run before publishing.}

% ============================================================================
\section{Lessons Learned}
% ============================================================================
\begin{enumerate}[noitemsep]
  \item \textbf{The market regime is the real variable.} No single strategy is
        universally right; mean reversion and momentum are two halves of one
        regime-conditional whole.
  \item \textbf{Negative results are load-bearing.} The failures of Phases 1--2
        are what \emph{justified} the momentum inversion and the regime work. We
        kept them in the story rather than hiding them.
  \item \textbf{Earn your complexity, then cut it.} We built the full ensemble,
        learned which part carried the edge, and simplified back to a gated
        momentum core --- more robust and less prone to overfitting.
  \item \textbf{For an unsupervised bot, risk plumbing \emph{is} the strategy.}
        Half the final code is circuit breakers, fill reconciliation, and error
        recovery, because a bot that blows up or hangs scores nothing.
\end{enumerate}

% ============================================================================
\appendix
\section{Note on Strategy 2}
\label{app:strat2}
The \code{strategy 2/} directory contains an independent module
(\code{strategy2\_hfm}) contributed by a teammate. It implements an academic
high-frequency-momentum signal: for several EMA pairs it computes a MACD-style
crossover $x=\mathrm{EMA}_{\text{short}}-\mathrm{EMA}_{\text{long}}$, normalises it
twice (by short- and long-horizon volatility) to obtain a $z$-score, and passes it
through a bounded nonlinear response $u = z\,e^{-z^2/4}/(\sqrt{2}\,e^{-1/2})$ that
\emph{attenuates} extreme readings. The components are averaged into a single
long/cash signal and backtested across $\sim$90 coins with hourly rebalancing,
next-bar execution, fees, and slippage. It is documented here only to delineate the
boundary of Strategy~1; it is not part of the research lineage described above.

\end{document}
